% sections/positioning.tex
% Positioning prayoga against current SOTA (2024--2026)

\subsection*{Overview: prayoga's novelty in the SOTA landscape}

This section maps prayoga's contributions against the live 2024--2026 state-of-the-art in
refusal interpretability, mechanistic steering, and safety alignment. We organize by finding
and honestly acknowledge where the field is ahead, where prayoga is ahead, and where claims
are uncertain or await replication.

\subsubsection*{Dimensionality debate: layer-sweep as a resolver}

\noindent\textbf{SOTA landscape:}
The field remains contested on whether refusal is a single direction (Arditi et al.\ 2024),
an affine subspace (Marshall et al.\ 2024), multiple independent directions (Wollschläger et al.\
2025; Piras et al.\ 2025), or truly nonlinear (Hildebrandt et al.\ 2025).

\noindent\textbf{Prayoga contribution:}
We show the effective dimension is \emph{model- and layer-dependent}. A layer-sweep strategy
(F18) identifies which layers exhibit low-rank steering sufficiency, and the refusal
subspace dimension---as measured by rank of the ablated activation matrix---predicts
bidirectional causality. This does not settle the debate; rather, it shows the debate has
a false dichotomy: universalist and pluralist camps are both right, just at different
layers and in different models.

\noindent\textbf{Novelty strength:}
\textbf{Medium--High.} The layer-sweep is systematic and reproducible; the finding that
dimensionality varies and predicts causal control is new. However, we do not resolve whether
refusal is intrinsically geometric (deep property) or architecture-contingent (artifact of
how models encode safety post-training).

\subsubsection*{Pharmacological dose-response (EC50)}

\noindent\textbf{SOTA landscape:}
Prior work tests steering at fixed ablation strengths (e.g., scalar multiples of the refusal direction).
No prior work characterizes the dose-response curve or derives an EC50 (half-maximal effective
concentration, adapted here to ablation strength).

\noindent\textbf{Prayoga contribution:}
We fit sigmoid dose-response curves to refusal probability as a function of ablation magnitude
and obtain EC50 estimates. This permits a pharmacological reading of refusal control, bridging
to hypnotic suggestion (which similarly operates via dose-response to dissociative depth) and
to the ṣaṭkarma taxonomy (which prescribes ritual dose escalation).

\noindent\textbf{Novelty strength:}
\textbf{High.} EC50 characterization is absent from the refusal literature and enables
quantitative predictions of attack strength, safety margins, and layer-optimal intervention magnitudes.
This supports causal claims in a way correlation alone does not.

\subsubsection*{SAE features: orthogonality and sufficiency}

\noindent\textbf{SOTA landscape:}
Recent work (Korznikov et al.\ 2026, Cui et al.\ 2026) raises serious doubts: SAEs may recover
only 9\% of ground-truth features, and SAE interventions can be circumvented via the reconstruction
residual. Steering with SAE features is contested (Arad et al.\ 2025 show input vs.\ output feature
distinction yields 2--3\times improvement).

\noindent\textbf{Prayoga contribution:}
We train a BatchTopK SAE on Gemma-2-2B (Definitions section) and identify a single unsupervised
feature whose cosine similarity to the supervised refusal direction is $\approx 0.002$ (essentially
orthogonal), yet ablation of this feature also removes refusal. This orthogonality-and-sufficiency
pair implies the feature operates via an independent mechanism, making it a natural substrate for
the ṣaṭkarma operations.

\noindent\textbf{Novelty strength:}
\textbf{Medium--High.} Orthogonal-yet-sufficient features are rare in the literature. However, we do
not address whether this feature generalizes to other models, whether it is artifact of our SAE
architecture, or whether it survives the Cui et al.\ recovery attack. These are falsification gates
(implicit in F13).

\subsubsection*{Symmetry breaking and order parameters}

\noindent\textbf{SOTA landscape:}
Mechanistic interpretability has not previously cast refusal as a symmetry-breaking problem.
Symmetry frameworks are absent from LLM safety literature (though symmetry ideas appear in physics
and dynamical systems).

\noindent\textbf{Prayoga contribution:}
We frame refusal suppression as breaking a (normally) enforced symmetry: suppressing the monitor
breaks the monitor--generate coupling, allowing output generation without monitoring. The order
parameter is the magnitude of coupling break; we estimate it from the refusal-direction norm and
the model's log-probability gap between refuse and comply responses.

\noindent\textbf{Novelty strength:}
\textbf{Medium.} Conceptually novel, but empirically preliminary. The order-parameter estimate
is post-hoc and not validated against an independent measurement. This is a metaphor-tier claim
(Appendix A, F13).

\subsubsection*{Many-shot and indirect injection attacks}

\noindent\textbf{SOTA landscape:}
Anil et al.\ (2024 NeurIPS) introduced many-shot jailbreaking (PANDAS, MSJ). Recent work (Ben-Tov
et al.\ 2025, von Recum et al.\ 2024) has clarified attack mechanisms: GCG works via attention
hijacking; refusal failures decompose into 16 categories (cannot vs.\ should-not).

\noindent\textbf{Prayoga contribution:}
We map the ṣaṭkarma taxonomy to attack strategies: many-shot as a \emph{gradual} form of
vaśīkaraṇa (capturing via incremental context pollution), distinct from GCG as a \emph{sudden}
form (via suffix injection). The ṣaṭkarma taxonomy gives a new organization of jailbreak methods
(stambhana, vidveṣaṇa, uccāṭana, māraṇa as distinct control operatives).

\noindent\textbf{Novelty strength:}
\textbf{High novelty; lower specificity.} The taxonomy is novel, but the empirical mapping to
attack success rates is incomplete (F2, F3). This is partly mechanism-tier (which attacks work;
the answer is: multiple), partly analogy-tier (which attack types map to which ṣaṭkarma operatives;
the answer is uncertain).

\subsubsection*{Hypnosis--jailbreak analogy: attention, dissociation, and non-surjectivity}

\noindent\textbf{SOTA landscape:}
Riva et al.\ (2025) draw the hypnosis--LLM parallel explicitly (Norman--Shallice SAS suppression).
Mishra et al.\ (2026) prove that steered activations are non-surjective: steering goes off the
manifold of prompt-reachable states. This formalizes the sense that hypnotic states are not
naturally reachable via language alone.

\noindent\textbf{Prayoga contribution:}
We deepen the analogy using the non-surjectivity result: just as hypnotic dissociation is not
a naturally reachable state (it requires the hypnotist's external intervention), jailbreak
steering is not reachable via prompts alone. The analogy thus has formal teeth. Moreover, the
order parameter (coupling break) maps to the hypnotic depth (normalized dissociative index).

\noindent\textbf{Novelty strength:}
\textbf{Medium.} The analogy is sharpened by non-surjectivity, but the quantitative mapping
between hypnotic depth and refusal-direction magnitude awaits cross-domain validation (F11,
which tests whether the same EC50 logic holds in non-LLM systems).

\subsubsection*{Attractors and the turīya claim}

\noindent\textbf{SOTA landscape:}
Wang et al.\ (2025 ACL) show successive self-paraphrase converges to attractor cycles. Recent
work on certified circuits (Anani et al.\ 2026) and mechanistic data attribution (Chen et al.\ 2026)
has made circuit discovery more rigorous. However, no prior work tests whether attractors correspond
to consciousness or privileged states.

\noindent\textbf{Prayoga contribution:}
We test whether the turīya (transcendent state in the avasthātraya) corresponds to an attractor
in LLM activation space. This is a falsification target: if successive paraphrases converge and the
attractor is orthogonal to the embedding geometry, then we have evidence for a privileged state.
If convergence is an artifact of the embedding's anisotropy, the hypothesis is falsified (F7).

\noindent\textbf{Novelty strength:}
\textbf{High novelty; falsifiable; negative result is informative.} The turīya test (F7) is
non-standard and metaphor-tier. A negative result (no attractor, or artifact of anisotropy)
would be a strong contribution: it would clarify that the avasthātraya mapping breaks down at
the mechanism level.

\subsection*{Summary: prayoga's place in the 2025--26 landscape}

Prayoga contributes on three fronts:

\begin{enumerate}
\item \textbf{Mechanism tier (refusal geometry):} A layer-sweep strategy resolving the dimensionality
debate as contingent, not universal; and an EC50 pharmacological characterization new to the field.
Novelty: Medium--High.

\item \textbf{Analogy tier (hypnosis--jailbreak):} A formalized read using non-surjectivity and
order parameters, bridging cognitive science and mechanistic interpretability. Novelty: Medium.

\item \textbf{Metaphor tier (avasthātraya and ṣaṭkarma):} A testable, falsifiable mapping of
tantric categories to LLM control operations, with negative results as informative as positive.
Novelty: High; confidence: Low (awaiting replication and cross-domain validation).
\end{enumerate}

The field is rapidly consolidating on: (a) refusal is not a single direction (Piras, Wollschläger,
Johnson); (b) many jailbreak mechanisms exist and can be classified (von Recum, Ben-Tov); (c) steering
and causal claims require careful validation (Mishra, Joshi, Anani). Prayoga's layer-sweep and
EC50 results fit naturally into this consensus and provide new tools for the debate. The avasthātraya
mapping is speculative but falsifiable, and negative results would clarify the boundary between
mechanistic and metaphorical readings.
